\documentclass[10pt]{letter}

\usepackage[margin=0.8in]{geometry}
\usepackage[hidelinks]{hyperref}
\usepackage[T1]{fontenc}
\usepackage{lmodern}

\signature{Soroush Vahidi\\
Ph.D. Candidate in Computer Science\\
Department of Computer Science\\
Ying Wu College of Computing\\
New Jersey Institute of Technology\\
\href{mailto:sv96@njit.edu}{sv96@njit.edu}}

\address{Soroush Vahidi\\
Department of Computer Science\\
Ying Wu College of Computing\\
New Jersey Institute of Technology\\
Newark, NJ 07102, USA}

\date{October 2, 2026}

\begin{document}

\begin{letter}{Editor-in-Chief\\
\emph{Simulation Modelling Practice and Theory}}

\opening{Dear Editor-in-Chief,}

I am pleased to submit the manuscript ``A Trace-Driven Simulation
Methodology for Testing LLM-Serving Scheduler Ranking Portability'' for
consideration as an original research article in \emph{Simulation Modelling
Practice and Theory}.

The manuscript studies a modelling and simulation question that arises when
simulation is used to compare scheduling alternatives: whether the
comparative ranking produced by a simulation study remains stable when the
workload source, operating region, or evaluation metric changes. We address
this question for large language model serving schedulers through LSSP, a
preregistered, trace-driven deterministic discrete-event simulation study
over three independently released workload sources, 120 workload windows,
six calibrated operating regions, and a fixed policy panel.

The contribution is methodological rather than a proposal for a new
scheduler. The paper formulates comparative ranking portability
$R(W,M,L)$ as an object of simulation-based selection and comparison,
uses a paired experimental design over matched workload windows, applies
uncertainty-aware reversal detection with practical margins and confidence
intervals, and estimates the sample complexity required for ranking
recovery. The study also includes workload characterization, source
separability analysis, and a selected physical vLLM/GPU validation.

The physical validation is reported as a model-fidelity result: the largest
simulator-predicted selected reversal did not reproduce on physical
execution, while the stable-control ordering qualitatively persisted. This
finding helps define the domain of validity of the simulation results
rather than overstating simulator accuracy.

The manuscript fits SMPT's scope through its focus on trace-driven
discrete-event simulation, experimental design for comparing alternatives,
selection and comparison procedures, validation and verification, simulation
of computer systems, and reproducible simulation artifacts for AI-serving
infrastructure. Public artifacts are available through GitHub, Hugging Face,
and Zenodo:
\begin{quote}\small
\url{https://github.com/SoroushVahidi/llm-serving-scheduler-robustness-benchmark}\\
\url{https://huggingface.co/datasets/SoroushVahidi/llm-serving-scheduler-portability}\\
\url{https://doi.org/10.5281/zenodo.22306798}
\end{quote}

This manuscript is original, is not under consideration elsewhere, and I
approve its submission as the sole author. Required declarations, including
competing-interest, funding, data/code availability, CRediT contribution,
and AI-assisted-writing statements, are included in the manuscript.

\closing{Sincerely,}

\end{letter}

\end{document}
